\begin{proposition}[Divisor twisting and embedded associated points]\hypertarget{tag-01GO}{}\label{cdp:st:E17}
\StatementRecord{cdp:st:E17}{proved}{01GO}{the definitions and constructions stated in the proof}

For an effective Cartier divisor $D$ on a compact complex
space $Z$ and a locally free sheaf $\mathcal V$, the spaces
$H^0(Z,\mathcal V(-kD))$, viewed inside $H^0(Z,\mathcal V)$,
stabilize for $k\gg0$ at
\[
 \{s\in H^0(Z,\mathcal V):\operatorname{supp}(s)\cap|D|
       =\varnothing\}.
\]
In particular they eventually vanish when $Z$ has no embedded
associated components and $D$ meets every irreducible component.
For the graph of a nonconstant meromorphic map from a smooth connected threefold
$Y$ to $\mathbb P^1$, every full fibre is isomorphic to a Cartier divisor in $Y$
and is Cohen--Macaulay with no embedded associated components.
\end{proposition}


\begin{proof}

\begin{enumerate}[label=\arabic*.,leftmargin=*]
\item Let $C=\mathbb P^1$, let $p,q$ be distinct points, and form the
square-zero complex space
\[
 |Z|=C,\qquad\mathcal O_Z=\mathcal O_C\oplus\mathbb C_p,
 \qquad(a,u)(b,v)=(ab,av+bu).
\]
Near $p$ its analytic local ring is
$\mathbb C\{z,e\}/(ze,e^2)$, and away from $p$ it is the
ordinary smooth curve. These charts glue through the quotient
by $e$, so this specifies a compact connected complex space.
It has one irreducible component, of positive dimension, and an
embedded associated point at $p$. The divisor $D=q$ is effective
Cartier: its local equation is a coordinate at $q$ and a unit
near $p$. For every integer $k$ there is the exact, split
description as an $\mathcal O_C$-module
\[
 \mathcal O_Z(-kD)=\mathcal O_C(-kq)\oplus\mathbb C_p.
\]
It gives every cohomology group:
\[
 h^0=\max(1-k,0)+1,\qquad h^1=\max(k-1,0),\qquad
 h^j=0\quad(j\geq2).
\]
In particular the nonzero section of the square-zero summand
survives for every $k\geq1$, contrary to
\citet[Proposition~1.1, p.~78]{CDP98}.

\item The descending spaces in the statement are
finite dimensional, hence stabilize. A section in their
intersection belongs to $\bigcap_k I_D^k\mathcal V_z=0$ at
every $z\in|D|$, by Krull intersection in the analytic local
ring. Its coherent support is therefore disjoint from $D$.
Conversely a section with support disjoint from $D$ vanishes
near $D$ and is a section of $\mathcal V(-kD)$ for every $k$.
When there are no
embedded associated components, every nonzero submodule of a
locally free sheaf has support containing an irreducible
component of $Z$. Such a support cannot avoid $D$ under the
stated intersection hypothesis.

\item Locally write the meromorphic map as $[a:b]$, where $a,b$ have
no common prime factor. Its graph $Z$ has equation
$av-bu=0$ in $Y\times\mathbb P^1$. At $c=[u_c:v_c]$, the maps
\[
 \{av_c-bu_c=0\}\longrightarrow Z_c,\qquad y\longmapsto(y,c),
\]
and the restriction of the projection to $Y$ are inverse holomorphic
maps of analytic spaces: both local rings are
$\mathcal O_Y/(av_c-bu_c)$. The local equation is a nonzero
element of the regular local ring of $Y$, hence a nonzerodivisor.
The quotient is Cohen--Macaulay and Gorenstein and has no embedded
associated components. This identifies the full analytic fibre,
including its nilpotents and multiplicities.

\end{enumerate}
\end{proof}